\documentclass[a4paper,12pt]{article}
%%% Set margins, indents, and linespread %%%

\usepackage[left=30mm, right=15mm, top=20mm, bottom=20mm]{geometry} % Package arguments define page margins
\pagestyle{plain} % Remove default upper header and leave only a page number at the bottom center of the page
\parindent=1.25cm % First line indent 1.25 cm, approximately equal to 5 characters (as required by GOST)
\usepackage{indentfirst} % Add indent to the first paragraph
%\linespread{1.3} % Linespread (the closest to 1.5 in Word) - here we use a command OnehalfSpacing* later


%%% Set language parameters and font %%%

\usepackage[russian]{babel}                         % English is the only language for the document
\usepackage{siunitx}
\babelfont{rm}{Times New Roman}    % TMR as the primary font
\babelfont{tt}{Courier New}    % cyrilic monofont 
\sisetup{output-decimal-marker = {,}}
\usepackage{fontenc}
\usepackage{fontspec} %ru in verbatim
\usepackage{verbatim}
%%% Set heading style for chapters, sections, and subsections %%%

% \usepackage{calc} % for indentation in the table of contents (before "c.")

\setsecnumdepth{subsection} % Numbers are only assigned to the first three levels, i.e., numbered are chapters, sections, and subsections
\renewcommand*{\chapterheadstart}{} % Redefine the command which adds spacing above the heading to prevent it from doing this
\renewcommand*{\printchaptername}{\ifbool{cent}{}{\hspace{1.25cm}}} % Redefine the command which prints the word 'Chapter' to prevent it from doing this
%\renewcommand*{\printchapternum}{} % Same for the chapter number - no, here we need to keep the number
\renewcommand*{\chapnumfont}{\normalfont\bfseries} % Change the chapter number font style: normal size, bold
\renewcommand*{\afterchapternum}{\hspace{1em}} % Change the delimiter between the chapter number and the chapter title
\renewcommand*{\printchaptertitle}{\normalfont\bfseries\ifbool{cent}{\centering}{}} % Change the chapter title font style: normal size, bold, centered, uppercase

\renewcommand*{\printtoctitle}[1]{%
  \begin{center}
    \normalfont\bfseries СОДЕРЖАНИЕ
  \end{center}
  \vspace{-1.3pt} \vspace{-\cftbeforechapterskip}
  \hspace{\linewidth} \hspace{-19pt} \textbf{с.}
} % unbelievable hack for indentation in the table of contents (before "c.")

\setlength{\beforechapskip}{1.3pt}
\setlength{\afterchapskip}{1.3pt}


\setbeforesecskip{1.3pt} % Set spacing above the section heading 
\setaftersecskip{1.3pt} % Same spacing below the section heading
\setsecheadstyle{\raggedright\hspace{1.25cm}\normalfont\bfseries} % Change the section heading font style: right aligned, normal size, bold
\setbeforesubsecskip{1.3pt} % Set spacing above the subsection heading
\setaftersubsecskip{1.3pt} % Same spacing below the subsection heading
\setsubsecheadstyle{\raggedright\normalfont\bfseries}  % Change the subsection heading font style: right aligned, normal size, bold



\setrmarg{2.55em plus1fil} % Disable hyphenation in toc
\setlength{\cftbeforechapterskip}{1.3pt} % This command removes vertical spacing before chapter titles - no, keep the spacing
% \renewcommand{\aftertoctitle}{\afterchaptertitle \vspace{-1.3pt} \vspace{-\cftbeforechapterskip} \begin{flushleft}\hfill с.\end{flushleft} 
% \vspace{-3.9pt}
% \vspace{-\cftbeforechapterskip}
% } % Adjust vertical spacing right after the heading 'Table of Contents' to make it the same as for other chapter headings
% \renewcommand{\aftertoctitle}{\\[0.5\baselineskip]\hfill\textnormal{с.}}
\renewcommand*{\chapternumberline}[1]{} % Redefine the command which prints the chapter number to prevent it from doing this - no, keep the number here
% \renewcommand*{\cftchapternumwidth}{1.5em} % Change the delimiter between the chapter number and the chapter title
\renewcommand*{\cftchapterfont}{\normalfont} % Chapter titles are printed with normal size, uppercase
\renewcommand*{\cftchapterpagefont}{\normalfont} % Page numbers are printed with normal size
\renewcommand*{\cftchapterdotsep}{\cftdotsep} % Create dot separators which go up to the page number for each entry in toc
\renewcommand*{\cftdotsep}{0} % Define the dot frequency
\renewcommand*{\cftchapterleader}{\cftdotfill{\cftchapterdotsep}} % Make the dots of standard shape (by default they are 'bold')
\maxtocdepth{section} % Only the first three levels are displayed in toc: chapters, sections, and subsections



%%% Alignment and hyphenation %%%

%% http://tex.stackexchange.com/questions/241343/what-is-the-meaning-of-fussy-sloppy-emergencystretch-tolerance-hbadness
%% http://www.latex-community.org/forum/viewtopic.php?p=70342#p70342
\tolerance 1414
\hbadness 1414
\emergencystretch 1.5em                             % If something goes wrong, try to adjust this one first
\hfuzz 0.3pt
\vfuzz \hfuzz
%\dbottom
%\sloppy                                            % No overfilling
\clubpenalty=10000                                  % Prevent page break after the first line in a paragraph
\widowpenalty=10000                                 % Prevent page break before the last line in a paragraph
\brokenpenalty=4991                                 % Prevent page break directly after hyphenation

%%% Set parameters for figures and tables %%%

\usepackage{graphicx, caption, subcaption} % Load packages for figures and captions
\graphicspath{{images/}} % Define the image folder
\captionsetup[figure]{font=normal, width=\textwidth, name=Рисунок, justification=centering, parskip=6pt} % Set the figure caption parameters: small font size (12pt), caption width is equal to the width of the main text, complete word "Figure", caption is centered
\captionsetup[subfigure]{font=normal} % Subfigure indices a), b), ... are also in small size (12pt) (by default they are even smaller)
\captionsetup[table]{singlelinecheck=false,font=normal,width=\textwidth,justification=justified, parskip=4pt, name=Таблица} % Set the table caption parameters: disable hyphenation, small font size (12pt), caption width is equal to the width of the main text, caption is justified
\captiondelim{ --- } % Set em dash as the delimiter between the caption name (i.e., the words 'Figure' and 'Table') and caption text
\setkeys{Gin}{width=\textwidth} % By defalut, the size of all figures will be adjusted to fit into the main text
%\setlength{\abovecaptionskip}{0pt} % Vertical space above the caption - no, keep the default
%\setlength{\belowcaptionskip}{0pt} % Vertical space below the caption - no, keep the default
\usepackage[section]{placeins} % All objects of the float type (figures/tables) do not go beyond the section where they are initialized

%%% Set parameters for references and hyperreferences %%%

\usepackage[unicode=ture]{hyperref}                               % Load the corresponding package
\hypersetup{
    colorlinks=true,                                % All refs and hyperrefs appear colored in text
    linktoc=all,                                    % Clickable colored refs for all entries in toc
    linktocpage=false,                               % Only the page numbers will be colored in toc, not the whole titles
    linkcolor=black,                                  % The color of refs and hyperrefs is red
    citecolor=black,                                 % The color of citations is red
    urlcolor = black
}

\usepackage{bookmark}
\bookmarksetup{
  numbered=true,    
  open=true,        
  depth=2            
}

%%% Set list parameters %%%

\usepackage{enumitem}                               % Load the package for the fine adjustment of list appearance
\renewcommand*{\labelitemi}{\normalfont{--}}        % Use en dash for unnumbered list items
\setlist{noitemsep, leftmargin=*,
%nolistsep
%,itemindent=\parindent
}                   % Remove vertical spacings between the items of the same level
\setlist[1]{labelindent=\parindent}                 % Item indent is equal to the paragraph indent
%\setlist[2]{leftmargin=\parindent}                  % Additional indent for the 2nd level
%\setlist[3]{leftmargin=\parindent}                  % And one more for the 3rd level

%%% Counters and object numbering %%%

\counterwithout{figure}{chapter}                    % Continuous numbering of figures throughout the document
\counterwithout{equation}{chapter}                  % Continuous numbering of math expressions throughout the document
\counterwithout{table}{chapter}                     % Continuous numbering of tables throughout the document

%%% Create bibliography using the biblatex package and biblatex-gost style processed by biber %%%

\usepackage{csquotes} % This package manages complex citation blocks (biblatex recommends to load it)
\usepackage[%
backend=biber,  % Processor
bibencoding=utf8,                                   % Encoding of the bib-file
sorting=none,                                       % Preference for sorting the reference list
style=gost-numeric,                                 % Citation style (GOST)
language=russian,                                   % Russian language for the reference list
sortcites=true,                                     % Several cites in square brackets will appear sorted
movenames=false,                                    % Do not move author names, they are always in the beginning of the citation entry
maxnames=5,                                         % Maximum number of authors for an entry
minnames=3,                                         % If there is more that 'max', cut the number of authors up to this
doi=false,                                          % Do now show DOI
isbn=false,    % Do not show ISBN, ISSN, ISRN
locallabelwidth=true,

]{biblatex}
\DeclareDelimFormat{bibinitdelim}{}                 % Remove the space between initials (Readdler J.R. instead of Readdler J. R.)
\addbibresource{bibl.bib}                           % Define the bibliography file



%%% Additional packages for enhanced functionality %%%

\usepackage{longtable,ltcaption}                    % Long tables
\usepackage{multirow,makecell}                      % Enriched formatting for tables
\usepackage{booktabs}                               % Yet another package for nice tables
\usepackage{soulutf8}                               % Support hyphenatable letterspacing, underline, and strikethrough
\usepackage[none]{hyphenat}                               % Enhanced hyphenation functionality
\usepackage{textcomp}                               % Support for 'complex' printed signs like some currency symbols, copyright, etc.
\usepackage[version=4]{mhchem}                      % Beautiful chemical equations
\usepackage{amsmath} % Enriched appearance of mathematical expressions 
\usepackage{amsthm}
\numberwithin{equation}{chapter}  % нумерация ур-ний вида 1.1, 1.2 (если chapter = 1)

%%% Subsequently add all document parts %%%

\setlength{\parskip}{0pt}


\usepackage{pdfpages}
\usepackage{float}
\usepackage{amssymb} % Для \mathbb{R}
\usepackage{tikz}
\usetikzlibrary{arrows.meta}
\usepackage[e]{esvect} %nice vec
\usepackage{mathtools} %rcases

\newtheorem{theorem}{Теорема}
\renewcommand\qedsymbol{$\blacksquare$}

\newcommand{\R}{\mathbb{R}}


\newbool{cent}

\hypersetup{
    colorlinks=true,
    linkcolor=black,
    urlcolor=blue
}

\begin{document}

\section*{Условие задачи}

\begin{table}[H]
    \centering
    \begin{tabular}{|c|c|c|c|c|c|c|}
        \hline
        $x$ & 2 & 4 & 6 & 8 & 10 & 12 \\
        \hline
        $y$ & 1,08 & 0,36 & 0,21 & 0,12 & 0,09 & 0,04 \\
        \hline
    \end{tabular}
    \caption{Исходные экспериментальные данные}
    \label{tab:data}
\end{table}

\section*{Задание}

\begin{enumerate}
    \item Описать исходные данные, привести примеры задач, в которых могут возникнуть подобные исходные данные, описать используемый метод решения задачи и обосновать его применение;
    \item Проанализировать исходные данные. Изобразить исходные данные графически при помощи библиотеки matplotlib, привести код для построения значений. Подобрать вид функции, наилучшим образом описывающие исходные данные, обосновать выбор именно этого вида функции;
    \item Решить задачу аналитически. Описать нахождение коэффициентов выбранной функции методом наименьших квадратов вручную, с построением промежуточных таблиц значений. Полученный результат изобразить графически;
    \item Решить задачу на языке Python с использованием библиотеки NumPy, привести код, привести результат решения;
    \item Решить задачу на языке Python с использованием библиотеки SciPy и Matplotlib;
    \item Сравнить решения, полученные вручную и с помощью Python, оценить точность;
    \item С помощью полученной модели предсказать значения y для $x = 4,5$ и $x = 5$.
\end{enumerate}

\newpage

\section{Решение задачи 1}

\subsection{Область применения}

Даны экспериментальные данные:

\begin{table}[H]
    \centering
    \begin{tabular}{|c|c|c|c|c|c|c|}
        \hline
        $x_i$ & 2 & 4 & 6 & 8 & 10 & 12 \\
        \hline
        $y_i$ & 1,08 & 0,36 & 0,21 & 0,12 & 0,09 & 0,04 \\
        \hline
    \end{tabular}
\end{table}

Подобные данные возникают при обработке результатов эксперимента, например, при исследовании зависимости физической величины $y$ от параметра $x$ (времени, температуры, нагрузки и т.д.). Как правило, значения $y_i$ содержат погрешности измерений.

Для аппроксимации зависимости используем метод наименьших квадратов. Конкретный вид функции будет выбран позже, а пока рассмотрим алгоритм для экспоненциальной функции.

Метод наименьших квадратов применяется для нахождения коэффициентов, минимизирующих сумму квадратов отклонений:

\[
S(\mathbf{a}) = \sum_{i=1}^{n} (y_i - f(x_i, \mathbf{a}))^2 \to \min_{\mathbf{a}}
\]

\subsection{Анализ исходных данных}

Рассмотрим исходные данные:

\[
x = 2, 4, 6, 8, 10, 12, \quad y = 1.08, 0.36, 0.21, 0.12, 0.09, 0.04.
\]

Для визуального анализа построим график экспериментальных точек.

\begin{figure}[H]
    \centering
    \includegraphics[width=0.7\textwidth]{1.png}
    \caption{Пример программы}
    \label{fig:raw_data}
\end{figure}

\begin{figure}[H]
    \centering
    \includegraphics[width=0.7\textwidth]{2.png}
    \caption{Визуализация исходных данных}
    \label{fig:raw_data}
\end{figure}

На основе визуального анализа графика можно сделать следующие выводы о характере зависимости:
\begin{itemize}
    \item функция монотонно убывает на всём рассматриваемом промежутке $x \in [2; 12]$;
    \item скорость убывания постепенно уменьшается;
    \item зависимость не является линейной.
\end{itemize}

Проверим гипотезу о линейности, вычислив первые разности:

\[
\Delta y_1 = 0.36 - 1.08 = -0.72, \quad \Delta y_2 = 0.21 - 0.36 = -0.15,
\]
\[
\Delta y_3 = 0.12 - 0.21 = -0.09, \quad \Delta y_4 = 0.09 - 0.12 = -0.03,
\]
\[
\Delta y_5 = 0.04 - 0.09 = -0.05.
\]

Значения $\Delta y_i$ не являются постоянными, что подтверждает нелинейный характер зависимости.

Учитывая, что функция убывает и имеет вогнутую форму, а также характерную для многих физических процессов форму, целесообразно рассмотреть экспоненциальную модель:

\[
y = a e^{bx}, \quad \text{где } b < 0.
\]

\subsection{Ручное нахождение коэффициентов МНК}

Для экспоненциальной модели $y = a e^{bx}$ выполним линеаризацию, прологарифмировав обе части:

\[
\ln y = \ln a + bx.
\]

Введём обозначения:

\[
Y = \ln y, \quad A = \ln a, \quad B = b.
\]

Получаем линейную зависимость:

\[
Y = A + Bx.
\]

Метод наименьших квадратов приводит к системе нормальных уравнений:

\[
\begin{cases}
n \cdot A + B \cdot \sum x_i = \sum Y_i, \\
A \cdot \sum x_i + B \cdot \sum x_i^2 = \sum x_i Y_i.
\end{cases}
\]

Для составления системы вычислим все необходимые суммы. Составим вспомогательную расчётную таблицу:

\begin{table}[H]
    \centering
    \begin{tabular}{|c|c|c|c|c|c|}
        \hline
        $i$ & $x_i$ & $y_i$ & $Y_i = \ln y_i$ & $x_i^2$ & $x_i Y_i$ \\
        \hline
        1 & 2 & 1,08 & 0,076961 & 4 & 0,153922 \\
        2 & 4 & 0,36 & -1,021651 & 16 & -4,086604 \\
        3 & 6 & 0,21 & -1,560648 & 36 & -9,363888 \\
        4 & 8 & 0,12 & -2,120264 & 64 & -16,962112 \\
        5 & 10 & 0,09 & -2,407946 & 100 & -24,079460 \\
        6 & 12 & 0,04 & -3,218876 & 144 & -38,626512 \\
        \hline
        \ Сумма & \ 42 & \ 190 & \ -10,252424& \ 364 & \ -92,964654 \\
        \hline
    \end{tabular}
    \caption{Таблица для расчёта коэффициентов}
    \label{tab:calc}
\end{table}

Вычислим необходимые суммы:
\[
n = 6, \quad \sum x_i = 42, \quad \sum Y_i = -10,252424, \quad \sum x_i^2 = 364, \quad \sum x_i Y_i = -92,964654.
\]

Решаем систему методом Крамера.

Вычислим определитель системы:
\[
\Delta = n \sum x_i^2 - \left( \sum x_i \right)^2 = 6 \cdot 364 - 42^2 = 2184 - 1764 = 420.
\]

Вычислим $\Delta_A$:
\[
\Delta_A = \sum Y_i \sum x_i^2 - \sum x_i \sum x_i Y_i = (-10,252424) \cdot 364 - 42 \cdot (-92,964654)  = 172,633132.
\]

Вычислим $\Delta_B$:
\[
\Delta_B = n \sum x_i Y_i - \sum x_i \sum Y_i = 6 \cdot (-92,964654) - 42 \cdot (-10,252424) -127,186116.
\]

Находим коэффициенты:
\[
A = \frac{\Delta_A}{\Delta} = \frac{172,633132}{420} = 0,411031,
\]
\[
B = \frac{\Delta_B}{\Delta} = \frac{-127,186116}{420} = -0,302824.
\]

Возвращаемся к исходным параметрам:
\[
a = e^A = e^{0,411031} = 1,5085,
\]
\[
b = B = -0,3028.
\]

Итоговая аппроксимирующая функция:

\[
\boxed{y = 1,5085 \cdot e^{-0,3028x}}
\]

\subsection{Графическое представление результатов}

\begin{figure}[H]
    \centering
    \includegraphics[width=0.7\textwidth]{3.png}
    \caption{Пример программы}
    \label{fig:raw_data}
\end{figure}

\begin{figure}[H]
    \centering
    \includegraphics[width=0.7\textwidth]{4.png}
    \caption{Экспоненциальная аппроксимация экспериментальных данных}
    \label{fig:analytic}
\end{figure}

\subsection{Решение с использованием библиотеки NumPy}

Для численного решения воспользуемся функцией \texttt{np.polyfit}, которая выполняет полиномиальную аппроксимацию. Для линеаризованных данных применим полином первой степени.

\begin{figure}[H]
    \centering
    \includegraphics[width=0.7\textwidth]{5.png}
    \caption{результат программы}
    \label{fig:raw_data}
\end{figure}

\begin{figure}[H]
    \centering
    \includegraphics[width=0.7\textwidth]{6.png}
    \caption{Аппроксимация с использованием NumPy}
    \label{fig:raw_data}
\end{figure}

\subsection{Решение с использованием библиотеки SciPy}

Библиотека SciPy предоставляет функцию \texttt{curve\_fit}, которая позволяет выполнять нелинейную аппроксимацию без предварительной линеаризации.

\begin{figure}[H]
    \centering
    \includegraphics[width=0.7\textwidth]{7.png}
    \caption{результат программы}
    \label{fig:raw_data}
\end{figure}

\begin{figure}[H]
    \centering
    \includegraphics[width=0.7\textwidth]{8.png}
    \caption{Аппроксимация с использованием NumPy}
    \label{fig:raw_data}
\end{figure}

\subsection{Сравнение решений и оценка точности}

Сопоставим коэффициенты экспоненциальной модели $y = a e^{bx}$, найденные тремя способами:

\begin{table}[H]
    \centering
    \begin{tabular}{|c|c|c|c|}
        \hline
        \ Коэффициент & \ Ручной расчёт & \ NumPy & \ SciPy \\
        \hline
        $a$ & 1,508500 & 1,508286 & 1,508700 \\
        $b$ & -0,302824 & -0,302824 & -0,302700 \\
        \hline
    \end{tabular}
    \caption{Сравнение коэффициентов}
    \label{tab:coeff_comp}
\end{table}

Все три метода дают практически идентичные результаты. Незначительные расхождения ($< 10^{-4}$) обусловлены погрешностями округления при ручных вычислениях и особенностями численных алгоритмов.

\subsection{Прогнозирование значений}

С помощью полученной модели предскажем значения $y$ для $x = 4,5$ и $x = 5$.

Выполним подстановку значений в уравнение модели (используем коэффициенты, полученные SciPy):

Для $x = 4,5$:
\[
\begin{aligned}
\hat{y}(4,5) &= 1,5087 \cdot e^{-0,3027 \cdot 4,5} \\
&= 1,5087 \cdot e^{-1,36215} \\
&= 1,5087 \cdot 0,2562 = 0,3866
\end{aligned}
\]

Для $x = 5,0$:
\[
\begin{aligned}
\hat{y}(5,0) &= 1,5087 \cdot e^{-0,3027 \cdot 5,0} \\
&= 1,5087 \cdot e^{-1,5135} \\
&= 1,5087 \cdot 0,2202 = 0,3323
\end{aligned}
\]

\begin{table}[H]
    \centering
    \begin{tabular}{|c|c|c|c|}
        \hline
        $x$ & \ Ручной расчёт & \ NumPy & \ SciPy \\
        \hline
        4,5 & 0,387 & 0,386 & 0,387 \\
        5,0 & 0,333 & 0,332 & 0,332 \\
        \hline
    \end{tabular}
    \caption{Прогнозные значения}
    \label{tab:predict}
\end{table}

\begin{figure}[H]
    \centering
    \includegraphics[width=0.7\textwidth]{9.png}
    \caption{Прогнозирование значений функции}
    \label{fig:predict}
\end{figure}

\begin{figure}[H]
    \centering
    \includegraphics[width=0.7\textwidth]{10.png}
    \caption{Прогнозирование значений функции}
    \label{fig:predict}
\end{figure}

\subsection{Выводы}

В ходе выполнения работы была построена экспоненциальная модель зависимости $y = a e^{bx}$ методом наименьших квадратов. Основные результаты:

\begin{enumerate}
    \item Получено уравнение регрессии:
    \[
    y = 1,5085 \cdot e^{-0,3028x}
    \]
    
    \item Коэффициенты, найденные аналитически (метод Крамера) и численно (NumPy, SciPy), совпадают с точностью до $10^{-4}$, что подтверждает корректность вычислений.
    
    \item Прогноз значений для $x = 4,5$ и $x = 5,0$ выполнен с оценкой точности; результаты могут быть использованы для интерполяции в диапазоне $x \in [2; 12]$.
\end{enumerate}

Метод наименьших квадратов позволил построить адекватную математическую модель экспериментальных данных. Сочетание аналитического расчёта и программной реализации обеспечивает как глубокое понимание метода, так и практическую эффективность решения.

\end{document}